\section*{动机与总览}

在本书中，我们将讨论 $(\infty,1)$-\trCategories{}的若干基础，以下简称%
\trInfinityCategories{}。\trInfinityCategories{}的例子由\trOrdinaryCategory{}%
以及拓扑（或单纯）\trEnrichment{}的\trCategories{}给出；此外，\trTopologicalSpace{}%
会产生\trInfinityGroupoid{}。事实上，相应的\trInfinityGroupoid{}只依赖于该%
\trTopologicalSpace{}的（弱）\trHomotopy{}型，而且任何\trInfinityGroupoid{}%
都以这种方式产生——这正是我们将在第 2 章看到的\trHomotopyHypothesis{}%
的内容。全体\trInfinityGroupoid{}的汇集本身就是一个很好的\trInfinityCategory{}，
它在某种意义上充当了集合的\trCategory{}在通常的\trCategoryTheory{}中所扮演的角色。

\trInfinityCategories{}自然地出现在\trHomotopyTheory{}与导出（代数）几何的诸多场合，
并且它们有一些非常好的特点：
\begin{enumerate}[label=(\arabic*)]
  \item 先前在\trModelCategory{}中不过是同伦\trLimit{}与同伦\trColimit{}之类的构造，
在\trInfinityCategories{}中变成了具有\trUniversalProperty{}的\trObject{}。
  \item 如果我们能证明某些\trObject{}具有\trUniversalProperty{}，那么就能构造出许多有趣的%
映射，其存在性无法仅凭考察相应\trObject{}的构造来预知（或者其存在依赖于手头恰好有%
某个特定的构造），例如从拓扑 $K$-理论（视为 $C^*$-代数上的\trFunctor{}）出发的自然映射，
或者\trTate{}对角（它在 TC 的新描述中至关重要）。
  \item 因此，\trInfinityCategory{}式的机制常常使人们能够对\trObject{}给出好的定义：%
例如，离散环的代数 $K$-理论如今可以完全按照\trQuillen{}与\trSegal{}设想的那样来定义。
不需要任何 $S^{-1}S$、$\mathcal{Q}$ 或 $S_\bullet$ 构造；毋宁说，它就是一个%
\trGroupCompletion{}，正如我们从 $K_0$ 中熟知的那样。这并不意味着这些构造没有用处——%
知道\trGroupCompletion{}可以用这些术语来描述，仍然是件好事。
  \item \trInfinityCategories{}的一个非常大的优点在于：由\trInfinityCategories{}构成的%
某些汇集会组装成一个很好的新\trInfinityCategory{}，而且人们由此可以在那一层次上%
进行许多范畴式的论证。例子包括\trFunctor{}的\trGoodwillie{}塔、%
\trGoerss{}--\trHopkins{}--\trMiller{}障碍塔、代数 $K$-理论作为非交换局部化不变量的%
\trUniversalProperty{}、其在\trHermitian{} $K$-理论中的变体、\trSpectrum{}的%
\trSmashProduct{}，等等。
  \item 借助\trInfinityCategories{}，把\trTopologicalSpace{}或\trScheme{}送到某些%
\trCategories{}（由\trSheaf{}或\trSheaf{}的复形构成）的指派满足\trDescent{}性质这一点，
在直观上变得清楚了；例如\trGrothendieck{}关于\trDerivedCategory{}的忠实平坦%
\trDescent{}定理。除其他作用外，这一事实还引出了对 $K$-理论\trDescent{}定理的现代处理%
与理解。类似地，借助它还可以定义\trMotivic{}$J$-同态
$K(S) \to \operatorname{Pic}(\operatorname{SH}(S))$，例如它被用于%
\trMotivic{}\trAtiyah{}对偶。
  \item 对有限\trCWComplex{}的广义\trPoincare{}对偶，有一种很好的处理方式，%
它基于研究这类空间上\trLocallyConstant{}\trSpectrum{}层，其精神与你或许从 $X$ 上的%
\trLocalSystem{}那里了解到的东西非常相似（事实上，\trLocalSystem{}无非是 $X$ 上%
\trLocallyConstant{}阿贝尔群\trSheaf{}）。由此可以得到关于闭流形、或更一般的%
\trPoincare{}对偶复形的\trPoincare{}对偶的一幅很好的图景；而且这一进路最清楚地%
说明了这类空间对哪类上同调论存在\trPoincare{}对偶，这一点非常有用，%
例如在\trSurgeryTheory{}中。
\end{enumerate}

上面的清单也许能让你相信\trInfinityCategories{}的理论在实践中是有用的。
然而，它也有一些缺点。依我之见，最大的缺点是下面这条（有意夸张的）：

\begin{warning*}
用手工构造出任何东西，实际上都是不可能的。
\end{warning*}

主要问题在于构造\trInfinityCategories{}之间的\trFunctor{}。我们将把%
\trInfinityCategories{}定义为某些种类的\trSimplicialSets{}，而一个\trFunctor{}将是%
\trSimplicialSets{}之间的映射。按照这一定义，\trFunctor{}不能这样来定义：先在%
\trObject{}与\trMorphism{}上指定对应，然后检验单位与合成是否得到保持。我们首先必须提供%
多得多的数据，这有时能用手工做到、但并非总能如此。因此，我们总是需要某种使我们得以%
构造此类\trFunctor{}的机制。

我们将要面对的问题的典型例子如下：
\begin{enumerate}[label=(\arabic*)]
  \item 设 $\Delta^1 \to J$ 是\trFreeStandingIsomorphism{} $J$ 中%
单个\trMorphism{}的包含。给定\trInfinityCategory{}中的一个\trMorphism{} $f$，
其\trClassifyingMap{}当且仅当 $f$ 是一个\trEquivalence{}时才经由 $\Delta^1 \to J$ 分解吗？
  \item 给定\trFunctor{}之间一个逐点为\trEquivalence{}的\trNaturalTransformation{}，
我们能找到该变换的逆吗？注意，取逆并不唯一，因此\trOrdinaryCategory{}中通常的证明%
在这里行不通。还要注意，作为一个\trNaturalTransformation{}，其含义多于一族具有某条%
性质（某些方块交换）的映射。
  \item 给定\trInfinityCategory{} $\mathcal{C}$ 中一族\trComposable{}\trMorphism{}%
$x_0 \to x_1 \to x_2 \to \dots$，我们能找到细化这些数据的\trFunctor{}%
$\mathbb{N} \to \mathcal{C}$ 吗？
  \item 给定一个\trFullyFaithful{}且\trEssentiallySurjective{}\trFunctor{}，
我们能找到它的逆吗？
\end{enumerate}

注意，所有这些问题在通常的\trCategoryTheory{}中答案都是肯定的。人们会期望它们在%
\trInfinityCategories{}中也成立，而本书的目标恰恰就是证明这类结果。原则上人们可以说，
任何来自通常\trCategoryTheory{}的合理事实都应当在\trInfinityCategories{}中有类似物，
但对「合理事实」究竟意味着什么，我们必须非常小心。下面是一些合理的（取决于你的背景）陈述：
\begin{enumerate}[label=(\arabic*)]
  \item 交换阿贝尔群/环的\trCategory{}是群/环的\trCategory{}的\trFullSubcategory{}。
  \item 遗忘\trFunctor{} $\mathcal{C}_{x/} \to \mathcal{C}$ 保持%
\trConnected{}\trColimit{}。
  \item 对常值\trFunctor{} $I \to * \xrightarrow{x} \mathcal{C}$ 取的\trColimit{}，
在 $I$ \trConnected{}时由 $x$ 给出。
\end{enumerate}

所有这些陈述在\trInfinityCategories{}的框架下都不成立。这与如下事实有关：%
\trInfinityCategories{}中的交换性不是一条性质，而是一份数据；特别地，映射可能保持、%
也可能不保持这样的数据。另一个事实更容易直观理解：如果把「\trConnected{}」换成%
「\trContractible{}」，那么类似的陈述将会成立。

本书最初的目标将是处理上述问题中的一部分。与此同时，我们将认识在该理论（的基础部分）%
中所用的那一类论证：\trSimplicialSets{}中的组合论证。在本书的最后部分，我们将讨论%
构造\trFunctor{}的另一个非常重要的工具：拉直--反拉直等价。我们不会证明它，但会讨论%
它的若干方面。随后，我们将用它来研究\trInfinityCategories{}之间的\trAdjunction{}，
并最终证明\trAdjointFunctorTheorem{}。
